% Section 7 -- Discussion. Source: ../manuscript.md section 7.
\section{Discussion}\label{sec:discussion}

The ladder is a branched accounting device, not a single cumulative
chain---and its two branches carry different kinds of content
(\sref{sec:ladder}): the measure branch the theorem content, the observer
branch protocol arithmetic. Across its dependency rows it recovers a substantial fraction of operational
Lorentzian geometry---dimension, timelike duration, radar decomposition,
Lorentz and Poincar\'e consistency, volume---from a causal order plus short,
explicit ingredient sets. It is equally a list of what each reconstruction
\emph{costs}: absolute scale costs a measure; a signed spatial coordinate
costs an orientation; a metric-not-merely-conformal statement costs a volume
element; Lorentzian structure is not bought by finite speed alone. Stating
both directions together is the point: it turns ``causal structure fixes the
conformal class, and a volume element fixes the remaining conformal factor''
from a slogan into a per-quantity operational ledger with measured error
behaviour.

This framing also clarifies what the program does \emph{not} show and sets up
the open question. Everything here is reconstruction inside known models: the
geometry is put in (by sprinkling from Minkowski or by a supplied protocol)
and recovered. It does not address whether geometry could \emph{emerge} from
an order that was not built from a geometry---the representability
question---which requires a validated discriminator and is taken up
separately. The \sref{sec:capstone} capstone sharpens the boundary from the
other side: the one piece of geometry that order itself carries---the
conformal class---is not merely present by theorem but legible to an order-only statistic at
finite density, exactly where the measure channel is frozen and only the
light cones move; how far that legibility extends before it fails is
unmeasured.

The Schwarzschild and mass-ladder stages add a second axis to that boundary.
Every verdict in the flat ladder and the plane-wave capstone is anchored
operationally---a margin taken from an exploration block, a threshold
declared before sizing. The \sref{sec:count} stage is the first
confrontation with a certified prediction instead, and what it buys is
deliberately narrow: not new physics, but the demonstration that an
operational finite-density instrument can be held to an interval-arithmetic
enclosure fixed before any data existed, at four masses, with the
near-degeneracy of the cross-rung volume comparison stated rather than
hidden. The provenance discipline that makes the confrontation
checkable---frozen decision rules, content-addressed evidence bundles,
contract tests that recompute every printed figure from its committed
artifact---is, we would argue, a contribution in its own right.

The usual scope statements are gathered here rather than in a separate
section. Results are controlled and mostly $(1{+}1)$D (dimension is checked
to 4D; the capstone and its extensions are $(3{+}1)$D). The constants are
convention-dependent (chain endpoint convention, null-inclusive causal
relation, Myrheim--Meyer normalization); we state each where it is used. The
spacelike proxy is exploratory. R5's local-measure evidence is one profile
family at one amplitude; the dose-response of the unweighted error floor to
the profile amplitude is not measured. Generalizing beyond type N is
partially opened rather than closed: the Schwarzschild extensions carry both
claim classes there (\sref{sec:schwarzschild}), and the certified oracle and
count ladder anchor the volume side (sections~\ref{sec:count}
and~\ref{sec:oracle}).

The natural next question---whether observer-relative distance \emph{order}
can be validated as recovering latent geometry, as opposed to being
reconstructed from a supplied one---is the subject of a companion paper (in
preparation) that builds a preregistered discriminator on this foundation,
measures its dose-response to geometry dilution and its dimension selection
in $(2{+}1)$D, and then carries it to orders produced by growth dynamics and
by an action-weighted 2D-order ensemble. At $N = 600$, random-start
post-burn-in configurations at $\beta = 2$ and $\beta = 8$ pass the frozen
instrument, while the $\beta = 32$ bipartite-start control is structurally
blocked (it fails at chain extraction, before any geometry gate). These are
not certified equilibrium draws, and the companion study did not establish
an equilibrium transition or finite-size scaling.
